\documentclass[12pt,fleqn]{amsart}

% Encoding
\usepackage[utf8]{inputenc}
\usepackage[T1]{fontenc}

% Page Layout
\usepackage{a4wide, booktabs}
\usepackage{scrextend}

% Maths Packages
\usepackage{amssymb,amsthm,latexsym}
\usepackage{mathrsfs}
\usepackage{dsfont}
\usepackage[mathscr]{euscript}
\usepackage{esint, enumitem, graphicx, float}
\usepackage{newtxtext,newtxmath} 
\usepackage{mathtools}
% Colour
\usepackage[usenames,dvipsnames,x11names]{xcolor}
\usepackage{listings}

% Define Mathematica syntax highlighting
\lstdefinelanguage{Mathematica}{
  morekeywords={ForAll, Resolve, Reals, Sqrt},
  sensitive=true,
  morecomment=[l]{(*},
  morecomment=[r]{*)},
  morestring=[b]{"}
}
\DeclareMathOperator{\ES}{ES}
% Web
\usepackage{url}

\usepackage{microtype}
\usepackage{hyperref}
\usepackage{pgfplots}
\pgfplotsset{compat=1.18}
\usepackage{tikz}

\renewcommand{\le}{\leqslant}
\renewcommand{\ge}{\geqslant}

\title[Dimension-dependent bounds for the SDIEP]{Dimension-dependent bounds for the SDIEP via phase optimisation and Paley-type constructions}
\author{Tomasz Kania}
\address{Institute of Mathematics, Czech Academy of Sciences, \v{Z}itn\'{a} 25, 115~67 Prague 1, Czech Republic \& Institute of Mathematics, Jagiellonian University, {\L}ojasiewicza 6, 30-348 Kraków, Poland}
\email{tomasz.marcin.kania@gmail.com}
\subjclass[2020]{15B51 (primary), 15A18, 15B48, 05B20, 05C50, 60J10 (secondary).}
\keywords{Symmetric doubly stochastic inverse eigenvalue problem, Sule\u{\i}manova spectrum, cycle basis, phase optimisation, coherence, Walsh--Hadamard matrix, Paley construction}
\dedicatory{In memoriam: Miroslav Fiedler (1926--2015),\\
on the occasion of the tenth anniversary of his passing}

\date{\today}
\thanks{IM CAS (RVO 67985840). }

\newtheorem{thm}{Theorem}
\newtheorem{prop}[thm]{Proposition}
\newtheorem{lem}[thm]{Lemma}
\newtheorem{cor}[thm]{Corollary}
\theoremstyle{remark}
\newtheorem{rem}[thm]{Remark}

\newcommand{\R}{\mathbb{R}}
\newcommand{\1}{\mathbf{1}}
\newcommand{\Lam}{\Lambda}

\begin{document}
\maketitle

\begin{abstract}
We refine the cycle-walk (Fourier) template of Gnacik and the author to quantify when a~$\delta$-Sule\u{\i}manova spectrum $(1,\lambda_2,\dots,\lambda_n)$ (with $\lambda_j\le 0$) is realised by a symmetric doubly stochastic matrix. For the canonical cycle basis we compute the \emph{exact} size-dependent threshold
\[
\delta_n \;=\; 1-\frac{1}{2\cos^2\!\Big(\frac{\pi}{4n}\rho(n)\Big)},
\quad
\rho(n)\in\{0,1,2,4\}\ \text{determined by } n\bmod 8,
\]
which improves $1/2$ if and only if $8\nmid n$; we also prove sharpness for that template. We then introduce an \emph{optimally phase-aligned} cycle basis which removes the `$8\mid n$' artefact and yields better sufficient bound
\[
\delta_n^{\rm (ph)} \;=\;
\begin{cases}
\displaystyle 1-\dfrac{1}{2\cos^2(\pi/n)},   & n\equiv 0\pmod{4},\\[2mm]
\displaystyle 1-\dfrac{1}{2\cos^2(\pi/2n)},  & n\equiv 2\pmod{4},\\[2mm]
\displaystyle 1-\dfrac{1}{2\cos^2(\pi/4n)},  & n\ \text{odd},
\end{cases}
\]
so that $\delta_n^{\rm (ph)}<\tfrac12$ for \emph{every} $n\ge3$ and $\delta_n^{\rm (ph)}=\delta_n$ unless $8\mid n$. Next, on abelian $2$-groups, the Walsh--Hadamard basis has coherence $M=1$ and hence suffices for \emph{all} Sule\u{\i}manova lists ($\delta=0$); the same conclusion holds in every Hadamard order (\emph{e.g.}, Paley families).
\end{abstract}

\section{Introduction}

A real square matrix is called doubly stochastic if it has non-negative entries and all of its row and column sums are equal to one. The Symmetric Doubly Stochastic Inverse Eigenvalue Problem (SDIEP) asks for the necessary and sufficient conditions under which a given list of real numbers $\sigma = (\lambda_1, \lambda_2, \dots, \lambda_n)$ can be the spectrum of an $n \times n$ symmetric doubly stochastic matrix \cite{GK20}. By the Perron--Frobenius theorem, any such matrix must have a Perron eigenvalue of $1$, corresponding to the eigenvector of all ones. This fixes $\lambda_1=1$ and, by Gershgorin's circle theorem, bounds all other eigenvalues to the interval $[-1, 1]$.\smallskip

A significant body of research focuses on the realisability of a \emph{Sule\u{\i}manova spectrum}, a list of the form $(1, \lambda_2, \dots, \lambda_n)$ where each $\lambda_j \le 0$ for $j \ge 2$ and the trace is non-negative, $1 + \sum_{j=2}^n \lambda_j \ge 0$. A powerful constructive method for tackling this problem is the Schur template: given a diagonal matrix $\Lambda = \operatorname{diag}(1, \lambda_2, \dots, \lambda_n)$, one constructs a candidate matrix $P(\Lambda) = Q\Lambda Q^\top$. If $Q$ is an orthogonal matrix whose first column is the Perron eigenvector $q_0 = \mathbf{1}/\sqrt{n}$, the resulting matrix $P(\Lambda)$ is automatically symmetric, doubly stochastic, and has the desired spectrum. 
\medskip

\noindent\textbf{Algorithmic template (Schur realisation).}

Given a basis $Q$ with $q_0=\1/\sqrt n$ and a list $\Lambda=\mathrm{diag}(1,\lambda_2,\ldots,\lambda_n)$:
\begin{enumerate}[leftmargin=*,itemsep=1pt]
\item Form $Y:=Q\Lambda$ \hfill (cost $O(n^2)$)
\item Form $P:=YQ^\top$ \hfill (dense GEMM, cost $O(n^3)$)
\item (Optional) Check entrywise nonnegativity: $\min_{k,\ell}P_{k\ell}\ge -\tau$ with tolerance $\tau$.
\end{enumerate}
For a chosen $Q$, Lemma~\ref{lem:coh} gives the sum-only sufficient condition $1+\sum_{j=2}^n\lambda_j>1-\tfrac{1}{M(Q)}$ that guarantees $P\ge0$. Lower $M(Q)$ \emph{tightens} this certificate.
\medskip

The entire problem is thus reduced to a single, challenging question: for a given spectrum $\Lambda$ and a chosen basis $Q$, what conditions guarantee that $P(\Lambda)$ is entrywise non-negative?\smallskip

The papers \cite{Perfect} and \cite{PM65} seem to be foundational for the problem. Recent work has approached this question from several angles. One approach provides general sufficient conditions that depend only on the sum of the non-Perron eigenvalues. A key result in this vein was provided by Gnacik and the author \cite{GK20}, who used the eigenbasis of the $n$-cycle graph for $Q$ to show that any Sule\u{\i}manova spectrum is realisable if $1 + \sum_{j=2}^n \lambda_j > 1/2$. It was noted, however, that this bound appeared to be sharp for this basis, particularly in dimensions where $n$ is a multiple of 8. A different approach, seen in the work of Sarkhoni \emph{et al}.~\cite{Sarkhoni25}, forgoes general sum-based conditions and instead derives necessary and sufficient conditions for highly structured spectra, such as $(1, 0, \dots, 0, \lambda_n)$, by constructing a bespoke orthogonal matrix $Q$ tailored to that specific list.\smallskip

In this paper, we unify and advance these perspectives. We begin by revisiting the cycle-walk basis from \cite{GK20}, replacing their baseline bound with a precise, dimension-dependent threshold that confirms the `$8 \mid n$' barrier is an artefact of the basis, not the problem itself. We then resolve this issue by introducing an optimally phase-aligned cycle basis that yields an improved bound, strictly better than $1/2$ for all dimensions $n \ge 3$. Finally, we generalise the underlying principle to a criterion based on basis \emph{coherence}, showing that bases derived from abelian 2-groups (Walsh--Hadamard matrices) achieve the ideal coherence, allowing for the realisation of \emph{all} Sule\u{\i}manova spectra with a strictly positive trace sum.\smallskip

Let $\sigma=(1,\lambda_2,\ldots,\lambda_n)$ with $\lambda_j\in[-1,0]$ for $j\ge 2$. Following Paparella \cite{Paparella16}, we call $\sigma$ a~\emph{Sule\u{\i}manova spectrum} and say it is \emph{$\delta$-Sule\u{\i}manova} if
\[
1+\sum_{j=2}^n\lambda_j>\delta .
\]
%The symmetric doubly stochastic inverse eigenvalue problem (SDIEP) asks when $\sigma$ is the spectrum of a symmetric, doubly stochastic matrix.

In \cite{GK20} a Markov-chain-inspired template was introduced: take $Q=[w_0,\ldots,w_{n-1}]$ to be a real orthonormal eigenbasis of the simple random walk on the $n$-cycle, so $w_0=\1/\sqrt n$ and, for $k\ge 1$,
\[
w_k(j)=\sqrt{\tfrac{2}{n}}\;\sin\!\Big(\tfrac{2\pi kj}{n}+\tfrac{\pi}{4}\Big)\quad (j=0,\ldots,n-1).
\]
Setting $P(\Lam)=Q\,\mathrm{diag}(1,\lambda_1,\ldots,\lambda_{n-1})\,Q^\top$ gives a symmetric matrix with unit row sums. Using $S_j(k):=\sin\!\big(\tfrac{2\pi kj}{n}+\tfrac{\pi}{4}\big)$ one has
\[
p_{kl}=\frac1n\!\left(1+2\sum_{j=1}^{n-1}\lambda_j\,S_j(k)\,S_j(l)\right).
\]
The crude estimate $S_j(k)S_j(l)\le 1$ and $\lambda_j\le 0$ yields the sufficient condition $1+\sum_{j=2}^n\lambda_j>\tfrac12$. It was also observed that when $8\mid n$ there are $j,k$ with $S_j(k)=1$, so this estimate cannot beat $\tfrac12$ for those sizes, and exhibited a $5\times 5$ spectrum at sum $0.48$ violating entrywise positivity. This suggests  trying other orthogonal choices of $Q$ to reduce the worst interaction $n\,|q_j(k)q_j(l)|$ (all from \cite{GK20}).

\medskip
In the present paper, we keep the Schur template $P(\Lam)=Q\Lam Q^\top$ and:
\begin{itemize}
\item compute \emph{exactly} the worst coefficient for the canonical cycle basis, obtaining a closed-form $\delta_n$ that improves $\tfrac12$ iff $8\nmid n$, with a matching sharpness example;
\item introduce an \emph{optimally phase-aligned} cycle basis (still orthogonal, same Perron column) giving the \emph{correct} piecewise bound $\delta_n^{\rm (ph)}<\tfrac12$ for every $n\ge3$, which strictly improves the $8\mid n$ case and is optimal within the cycle family;
\item prove a general `coherence $\Rightarrow$ sufficient $\delta$' lemma and apply it to Walsh--Hadamard bases (abelian $2$-groups) and Hadamard orders (via Paley), obtaining \emph{$\delta=0$ in those orders only}. Here `$\delta=0$' means the \emph{strict} condition $1+\sum\lambda_j>0$ suffices; this does not contradict obstructions on the boundary $1+\sum\lambda_j=0$ (\emph{e.g.}, $(1,0,\dots,0,-1)$ in odd dimensions).
\end{itemize}

\section{The cycle-walk basis}\label{sec:cycle}
From now on we reindex the list as $\sigma=(\lambda_0,\lambda_1,\ldots,\lambda_{n-1})$ with $\lambda_0=1$ and $\lambda_j\in[-1,0]$ for $1\le j\le n-1$, and use the cycle eigenbasis $Q=[w_0,\dots,w_{n-1}]$ described above.

\begin{lem}\label{lem:entrybound}
Let $Q$ and $\Lam$ be as above. Then
\begin{equation}\label{eq:entry}
p_{kl}
=\frac1n\!\left(1+2\sum_{j=1}^{n-1}\lambda_j\,S_j(k)\,S_j(l)\right),
\quad
S_j(k):=\sin\!\Big(\tfrac{2\pi kj}{n}+\tfrac{\pi}{4}\Big).
\end{equation}
If $\lambda_j\le 0$, writing
\[
c_n:=\max_{1\le j\le n-1}\ \max_{0\le k,l\le n-1} S_j(k)\,S_j(l),
\]
we have the bound
\begin{equation}\label{eq:lowerbd}
p_{kl}\ \ge\ \frac1n\Bigl(1+2c_n\sum_{j=1}^{n-1}\lambda_j\Bigr)\qquad(0 \le k, l \le n-1).
\end{equation}
\end{lem}

\begin{proof}
Orthogonality and $w_0=\1/\sqrt n$ give $P(\Lam)\1=\1$. Expanding $Q\Lam Q^\top$ yields \eqref{eq:entry}. Because $\lambda_j\le0$ and $S_j(k)S_j(l)\le c_n$ for each $j$, multiplying the latter inequality by $\lambda_j$ \emph{reverses} it, hence
\(
\sum_{j=1}^{n-1}\lambda_j S_j(k)S_j(l)\ \ge\ c_n\sum_{j=1}^{n-1}\lambda_j,
\)
and \eqref{eq:lowerbd} follows.
\end{proof}

\begin{lem}[Exact value of $c_n$ and the mod-$8$ mechanism]\label{lem:cn}
Define
\[
\Delta_n \;:=\; \min_{m,t\in\mathbb Z}\left|\frac{2\pi m}{n}-\frac{(2t+1)\pi}{4}\right|.
\]
Then
\[
c_n=\cos^2\Delta_n
\quad\text{and}\quad
\Delta_n=\frac{\pi}{4n}\,\rho(n),
\qquad
\rho(n)=
\begin{cases}
0,& n\equiv0\pmod{8},\\
1,& n\equiv1,3,5,7\pmod{8},\\
2,& n\equiv2,6\pmod{8},\\
4,& n\equiv4\pmod{8}.
\end{cases}
\]
Moreover, the maximum in $c_n$ is attained with $k=l$.
\end{lem}

\begin{proof}
For fixed $j$ with $\gcd(j,n)=1$, the multiset $\{2\pi kj/n+\pi/4:\ 0\le k\le n-1\}$ is a permutation of the grid $\{2\pi m/n+\pi/4:\ 0\le m\le n-1\}$. Hence
\[
\max_k|S_j(k)|=\max_m\Big|\sin\Big(\frac{2\pi m}{n}+\frac{\pi}{4}\Big)\Big|
=\sin\!\Big(\frac{\pi}{2}-\Delta_n\Big)=\cos\Delta_n,
\]
where $\Delta_n$ is the minimal angular distance from this grid to the odd half-angles. Therefore $c_n=(\max_{j,k}|S_j(k)|)^2=\cos^2\Delta_n$, and the maximum occurs already for $k=l$ at a peak. As
\[
\frac{2\pi m}{n}-\frac{(2t+1)\pi}{4}=\frac{\pi}{4n}\big(8m-(2t+1)n\big)
\]
implies $\Delta_n=\tfrac{\pi}{4n}\min_{m,t}|\,8m-(2t+1)n\,|$. We now analyse residues of $n$ modulo $8$.
\begin{itemize}
\item {$n\equiv0\pmod 8$.}
Then $(2t+1)n$ is divisible by $8$ for every $t$, so choosing $m=(2t+1)n/8$ gives the minimum $0$. Thus $\rho(n)=0$.

\item {$n$ odd.}
Here $\gcd(n,8)=1$. By Bézout's theorem, there exist integers $m,t$ with $8m-(2t+1)n=1$. Hence, the minimum positive value is $1$ and $\rho(n)=1$.

\item {$n\equiv 2,6\pmod 8$.}
Write $n=2k$ with $k$ odd. Then
\[
8m-(2t+1)n = 2\bigl(4m-(2t+1)k\bigr)
\]
is always even. Moreover $4m-(2t+1)k$ is always odd, since $4m$ is even and $(2t+1)k$ is odd, so $8m-(2t+1)n$ is never divisible by $4$. Hence the minimum positive value is at least $2$. Conversely, since $\gcd(4,k)=1$, there exist integers $m,r$ with $4m-rk=1$. Reducing modulo $2$ shows that $r$ is odd, say $r=2t+1$, and then
\[
8m-(2t+1)n = 2.
\]
Thus $\rho(n)=2$.

\item {$n\equiv4\pmod 8$.}
Write $n=4k$ with $k$ odd. Then
\[
8m-(2t+1)n = 4\bigl(2m-(2t+1)k\bigr)
\]
is always divisible by $4$. Since $2m-(2t+1)k$ is odd, it is never divisible by $2$, so $8m-(2t+1)n$ is never divisible by $8$. Hence the minimum positive value is at least $4$. Conversely, since $\gcd(2,k)=1$, there exist integers $m,r$ with $2m-rk=1$. Again $r$ must be odd, say $r=2t+1$, and then
\[
8m-(2t+1)n = 4.
\]
Thus $\rho(n)=4$.
\end{itemize}
\end{proof}

\begin{thm}[Dimension-dependent bound for the cycle basis]\label{thm:cycle}
With $\Delta_n$ as above, set
\[
\delta_n:=1-\frac{1}{2\cos^2\Delta_n}
\ =\ \frac12-\frac12\tan^2\Delta_n.
\]
If $1+\sum_{j=2}^n\lambda_j>\delta_n$, then $P(\Lam)=Q\Lam Q^\top$ is symmetric and doubly stochastic. In particular $\delta_n<\tfrac12$ if and only if $8\nmid n$, whereas $\delta_n=\tfrac12$ when $8\mid n$.
\end{thm}

\begin{proof}
By Lemma~\ref{lem:entrybound}, for every $k,l$,
\[
p_{kl}\ \ge\ \frac1n\Bigl(1+2c_n\sum_{j=1}^{n-1}\lambda_j\Bigr),
\qquad c_n=\cos^2\Delta_n.
\]
by Lemma~\ref{lem:cn}. Hence $p_{kl}\ge0$ whenever $\sum_{j=1}^{n-1}\lambda_j>-1/(2c_n)$, i.e.,\ when
\[
1+\sum_{j=2}^n\lambda_j\ >\ 1-\frac{1}{2\cos^2\Delta_n}\ =\ \delta_n,
\]
which establishes entrywise nonnegativity under the stated hypothesis.

\smallskip
Since $Q$ is orthogonal and $\Lambda$ is (real and) diagonal (hence $\Lambda^\top=\Lambda$),
\[
P(\Lambda)^\top=(Q\Lambda Q^\top)^\top=(Q^\top)^\top\Lambda^\top Q^\top=Q\Lambda Q^\top=P(\Lambda),
\]
so $P(\Lambda)$ is symmetric.

\smallskip
Let $q_0,\dots,q_{n-1}$ be the columns of $Q$ and let $e_0=(1,0,\dots,0)^\top$. Since $q_0=\1/\sqrt n$ and $q_j\perp q_0$ for $j\ge1$,
\[
Q^\top\1=\big(\langle q_0,\1\rangle,\dots,\langle q_{n-1},\1\rangle\big)^\top=(\sqrt n,0,\dots,0)^\top=\sqrt n\,e_0.
\]
Therefore
\[
P(\Lambda)\,\1
=Q\Lambda Q^\top\1
=Q\Lambda(\sqrt n\,e_0)
=\sqrt n\,Q(\Lambda e_0)
=\sqrt n\,Q e_0
=\sqrt n\,q_0
=\sqrt n\,\frac{\1}{\sqrt n}
=\1.
\]
Thus every row sum equals $1$. Finally,
\[
\1^\top P(\Lambda)
=\1^\top Q\Lambda Q^\top
=(Q^\top\1)^\top \Lambda Q^\top
=(\sqrt n\,e_0)^\top \Lambda Q^\top
\]
\[
=\sqrt n\,e_0^\top \Lambda Q^\top
=\sqrt n\,e_0^\top Q^\top
=(\sqrt n\,Q e_0)^\top
=\1^\top,
\]
where we used that $\Lambda e_0=e_0$ because the first diagonal entry of $\Lambda$ is $\lambda_0=1$. Hence, every column sum also equals $1$.

Combining the three parts shows that $P(\Lambda)$ is symmetric and doubly stochastic, whenever $1+\sum_{j=2}^n\lambda_j>\delta_n$. The stated equivalence for $\delta_n<\tfrac12$ versus $8\nmid n$ follows from Lemma~\ref{lem:cn}, which gives $c_n=\cos^2\Delta_n=1$ iff $8\mid n$.
\end{proof}


\begin{prop}[Sharpness for the cycle basis]\label{prop:sharp}
Fix $n$ and choose indices realising 
\[
    c_n=\max_{j,k,l}S_j(k)S_j(l).
\]
If $\lambda_{j^{\ast}}=t<0$ and $\lambda_j=0$ for $j\ne j^{\ast}$, then
\[
p_{k^{\ast}k^{\ast}}=\frac1n\bigl(1+2t\,\cos^2\Delta_n\bigr),
\]
which is non-negative if and only if $t\ge -1/(2\cos^2\Delta_n)$, i.e.\ $1+\sum_{j=2}^n\lambda_j\ge\delta_n$. Thus Theorem~\ref{thm:cycle} is optimal for this $Q$ among conditions depending only on $\sum_{j\ge 2}\lambda_j$.
\end{prop}

\begin{rem}We see that divisibility by $8$ is the exact barrier. Indeed, by Lemma~\ref{lem:cn}, the grid $\{2\pi m/n+\pi/4\}$ hits an odd half-angle iff $8\mid n$, and then $c_n=1$ so the factor $\tfrac12$ is unavoidable for this $Q$. If $8\nmid n$, the nearest grid point sits at angular distance $\Delta_n=\pi\rho(n)/(4n)$ from $\pi/2$, giving $c_n=\cos^2\Delta_n<1$ and hence $\delta_n<1/2$.
\end{rem}

\section{Phase-optimised cycle bases}

The eigenbasis of the $n$-cycle graph is not unique. For each $j \in \{1, \dots, \lfloor(n-1)/2\rfloor\}$, the corresponding eigenspace $E_j$ is two-dimensional, spanned by a sine and cosine vector. The canonical basis analysed in Section~\ref{sec:cycle} represents just one choice of basis within these spaces, a~choice that, for dimensions where $8 \mid n$, leads to an unavoidable coherence peak and a bound of $\delta_n=1/2$. This section demonstrates that by rotating the basis in each eigenspace via an optimally chosen phase angle $\phi$, we can construct a new orthogonal matrix $Q^{\rm (ph)}$. This `phase-optimised` basis is designed to minimise the maximum magnitude of its entries, effectively smoothing the coherence peaks and removing the ``$8 \mid n$'' artefact entirely. More specifically, for each $j\in\{1,\dots,\lfloor(n-1)/2\rfloor\}$ the cycle eigenspace is two-dimensional,
$E_j=\operatorname{span}\{u_j,v_j\}$ with
\(
u_j(k)=\sqrt{\tfrac2n}\cos(\tfrac{2\pi jk}{n}),
\ v_j(k)=\sqrt{\tfrac2n}\sin(\tfrac{2\pi jk}{n}) .
\)
Any orthonormal pair in $E_j$ is of the form
\[
w^{(\phi)}_j(k)=\sqrt{\tfrac{2}{n}}\sin\!\Big(\tfrac{2\pi jk}{n}+\phi\Big),\qquad
\tilde w^{(\phi)}_j(k)=\sqrt{\tfrac{2}{n}}\cos\!\Big(\tfrac{2\pi jk}{n}+\phi\Big),
\]
for some phase $\phi$. When $n$ is even and $j=n/2$, the eigenspace is one-dimensional; choosing the phase so that $|\sin(\phi)|=1/\sqrt2$ yields a unit vector with $|w_{n/2}(k)|\equiv 1/\sqrt n$, hence $n\,|q_{n/2}(k)q_{n/2}(l)|\equiv 1$.

We must control \emph{both} columns in each $E_j$ simultaneously. Write $n':=n/\gcd(j,n)$; then $\{\tfrac{2\pi jk}{n}\}$ is a translate of the grid with step $2\pi/n'$.

\begin{lem}[Two-phase covering radius on a grid]\label{lem:twophase}
Let $n':=n/\gcd(j,n)$ and set $S:=2\pi/n'$ and $L:=\pi/2$. All angles are considered reduced modulo $L$. Define
\[
\Delta^\star(n')\ :=\
\begin{cases}
\pi/n',     & n'\equiv 0\ (\mathrm{mod}\ 4),\\[1mm]
\pi/(2n'),  & n'\equiv 2\ (\mathrm{mod}\ 4),\\[1mm]
\pi/(4n'),  & n'\ \mathrm{odd}.
\end{cases}
\]
For $w^{(\phi)}_j(k)=\sqrt{\tfrac{2}{n}}\sin(\tfrac{2\pi jk}{n}+\phi)$ and
$\tilde w^{(\phi)}_j(k)=\sqrt{\tfrac{2}{n}}\cos(\tfrac{2\pi jk}{n}+\phi)$ one has
\[
\min_{\phi}\ \max\!\Big\{\,\|w^{(\phi)}_j\|_\infty,\ \|\tilde w^{(\phi)}_j\|_\infty\,\Big\}
\ =\ \sqrt{\tfrac2n}\ \cos\big(\Delta^\star(n')\big),
\]
and consequently
\[
\min_{\phi}\ \sup_{k,l}\ n\,\max\{|w^{(\phi)}_j(k)w^{(\phi)}_j(l)|,\ |\tilde w^{(\phi)}_j(k)\tilde w^{(\phi)}_j(l)|\}
\ =\ 2\cos^2\big(\Delta^\star(n')\big).
\]
\end{lem}

\begin{proof}
Let $A:=\{mL:\ m\in\mathbb Z\}$ with $L=\pi/2$. For any $x\in\mathbb R$,
\begin{equation}\label{eq:max-sin-cos}
\max\{|\sin x|,|\cos x|\}=\cos\big(\operatorname{dist}(x,A)\big).
\end{equation}
Indeed, reduce $x$ modulo $L$ to $y\in[0,L]$. Then
\[
|\cos x|=|\cos y|,\qquad |\sin x|=|\sin y|=|\cos(L-y)|.
\]
Hence $\max\{|\sin x|,|\cos x|\}=\max\{|\cos y|,|\cos(L-y)|\}=\cos\big(\min\{y,L-y\}\big)$.
However $\min\{y,L-y\}=\operatorname{dist}(x,A)$ by construction, giving \eqref{eq:max-sin-cos}.

\smallskip

Fix $j$ and write $n':=n/\gcd(j,n)$. The set of phases
\[
\Theta_j:=\Big\{\tfrac{2\pi jk}{n}\,:\ k=0,1,\dots,n-1\Big\}
\]
is a union of $\gcd(j,n)$ translates of the uniform grid $\{2\pi m/n':\ m=0,\dots,n'-1\}$. Since a~translation does not affect sup norms, it suffices to study the latter uniform grid. After adding a~free phase $\phi$ and reducing modulo $L$ we obtain
\[
X_\phi:=\Big\{\tfrac{2\pi m}{n'}+\phi\ \bmod L:\ m=0,\dots,n'-1\Big\}\ \subset \ \mathbb R/L\mathbb Z.
\]
This set is a coset (translate) of the cyclic subgroup $G:=\langle \overline{S}\rangle$ of the circle $\mathbb R/L\mathbb Z$, generated by the step $\overline{S}$ which is the class of $S:=2\pi/n'$ modulo $L$.

\smallskip

The order $r:=|G|$ is the least positive integer such that $r\overline{S}=\overline{0}$, i.e.
\[
rS\in L\mathbb Z \quad\Longleftrightarrow\quad \frac{2\pi r}{n'}=\frac{\pi}{2}\,q \ \text{ for some } q\in\mathbb Z
\quad\Longleftrightarrow\quad \frac{4r}{n'}=q\in\mathbb Z.
\]
Thus $r$ is the least positive integer with $n' \mid 4r$, which is
\[
r=\frac{n'}{\gcd(n',4)}.
\]
Consequently, the $r$ points of $G$ are equally spaced on the circle $\mathbb R/L\mathbb Z$ with spacing
\begin{equation}\label{eq:seff}
s_{\mathrm{eff}}=\frac{L}{r}=\frac{\pi}{2}\cdot\frac{\gcd(n',4)}{n'}
=
\begin{cases}
2\pi/n',   & n'\equiv 0\ (\mathrm{mod}\ 4),\\[1mm]
\pi/n',    & n'\equiv 2\ (\mathrm{mod}\ 4),\\[1mm]
\pi/(2n'), & n'\ \text{odd}.
\end{cases}
\end{equation}

\smallskip

For a fixed translate $X_\phi=\overline{\phi}+G$, define the nearest-point distance
\[
d(\phi):=\min_{x\in X_\phi}\operatorname{dist}(x,\overline{0}).
\]
Since $X_\phi$ consists of $r$ equally spaced points with gap $s_{\mathrm{eff}}=L/r$, the circle $\mathbb R/L\mathbb Z$ is partitioned into $r$ arcs of length $s_{\mathrm{eff}}$ whose endpoints are the points of $X_\phi$. The quantity $d(\phi)$ is the distance from $\overline{0}$ to the nearest endpoint of the arc containing $\overline{0}$, so
\[
d(\phi)\le \frac{s_{\mathrm{eff}}}{2},
\]
with equality exactly when $\overline{0}$ is the midpoint of such an arc. Therefore
\[
\max_{\phi} d(\phi)=\frac{s_{\mathrm{eff}}}{2}.
\]
For instance, if $x_0\in G$ is arbitrary, choosing $\phi$ so that $\overline{\phi}+x_0=\overline{s_{\mathrm{eff}}/2}$ attains this maximum.

\smallskip

For any $\phi$ we have, by \eqref{eq:max-sin-cos},
\[
\max\Big\{\|w^{(\phi)}_j\|_\infty,\ \|\tilde w^{(\phi)}_j\|_\infty\Big\}
=\sqrt{\tfrac{2}{n}}\ \max_{x\in X_\phi}\max\{|\sin x|,|\cos x|\}
=\sqrt{\tfrac{2}{n}}\ \max_{x\in X_\phi}\cos\big(\operatorname{dist}(x,\overline{0})\big).
\]
Since $\cos$ is decreasing on $[0,\pi/2]$, the inner maximum equals $\cos(d(\phi))$. Hence
\[
\min_{\phi}\ \max\Big\{\|w^{(\phi)}_j\|_\infty,\ \|\tilde w^{(\phi)}_j\|_\infty\Big\}
=\sqrt{\tfrac{2}{n}}\ \cos\!\Big(\max_{\phi} d(\phi)\Big)
=\sqrt{\tfrac{2}{n}}\ \cos\!\Big(\frac{s_{\mathrm{eff}}}{2}\Big).
\]
Combining this with \eqref{eq:seff} gives the stated value $\sqrt{2/n}\,\cos(\Delta^\star(n'))$.

Finally, for either column in the pair, we obtain
\[
\min_{\phi}\ \sup_{k,l}\ n\,|q(k)q(l)|
\ \le\ \min_{\phi}\ n\,\|q\|_\infty^2
\ \le\ n\cdot \Big(\sqrt{\tfrac{2}{n}}\cos(\Delta^\star(n'))\Big)^2
=2\cos^2(\Delta^\star(n')).
\]
This upper bound is tight because, at a minimising shift, there is some $k$ with
\[
\operatorname{dist}\!\Big(\frac{2\pi jk}{n}+\phi,\ A\Big)=d(\phi)=\frac{s_{\mathrm{eff}}}{2},
\]
hence for one of the two vectors we have
\[
|q(k)|=\sqrt{\tfrac{2}{n}}\cos(\Delta^\star(n')).
\]
Taking $l=k$ shows that the value $2\cos^2(\Delta^\star(n'))$ is attained.
\end{proof}


Because $\Delta^\star(n')$ decreases with $n'$, the worst case is $\gcd(j,n)=1$, \emph{i.e.},\ $n'=n$.

\begin{thm}[Phase-optimised coherence]\label{thm:phase-correct}
There exists an orthogonal matrix $Q^{\rm (ph)}$ with first column $\1/\sqrt n$ such that
\[
M\big(Q^{\rm (ph)}\big)\ =\
\begin{cases}
2\cos^2\!\big(\dfrac{\pi}{n}\big),   & n\equiv 0\pmod{4},\\[2mm]
2\cos^2\!\big(\dfrac{\pi}{2n}\big),  & n\equiv 2\pmod{4},\\[2mm]
2\cos^2\!\big(\dfrac{\pi}{4n}\big),  & n\ \mathrm{odd},
\end{cases}
\]
and this value is optimal within the cycle family.
Consequently,
\[
\delta^{\mathrm{(ph)}}_n \;:=\;
1-\frac{1}{ M(Q^{\rm (ph)})}
\ =\
\begin{cases}
\displaystyle 1-\dfrac{1}{2\cos^2(\pi/n)},   & n\equiv 0\pmod{4},\\[2mm]
\displaystyle 1-\dfrac{1}{2\cos^2(\pi/2n)},  & n\equiv 2\pmod{4},\\[2mm]
\displaystyle 1-\dfrac{1}{2\cos^2(\pi/4n)},  & n\ \text{odd},
\end{cases}
\]
and $0\le \delta^{\mathrm{(ph)}}_n<\tfrac12$ for every $n\ge 3$. Moreover, $\delta^{\mathrm{(ph)}}_n=\delta_n$ unless $8\mid n$, and in all cases $\delta^{\mathrm{(ph)}}_n \nearrow \tfrac12$ as $n\to\infty$.
\end{thm}

\begin{proof}
To show the upper bound, for each $j\in\{1,\dots,\lfloor (n-1)/2\rfloor\}$ let us choose the explicit phase
\[
\phi_j:=\Delta^\star\!\bigl(n/\gcd(j,n)\bigr),
\]
which attains the minimum in Lemma~\ref{lem:twophase} for $n'=n/\gcd(j,n)$, and replace the canonical pair $\{u_j,v_j\}$ by $\{w_j^{(\phi_j)},\tilde w_j^{(\phi_j)}\}$. This rotation within $E_j$ preserves orthonormality and keeps the Perron column $w_0=\1/\sqrt n$ unchanged; doing this independently for each $j$ yields an orthogonal matrix $Q^{\rm (ph)}$.

By Lemma~\ref{lem:twophase}, for each such $j$,
\[
\max\big\{\|w_j^{(\phi_j)}\|_\infty,\ \|\tilde w_j^{(\phi_j)}\|_\infty\big\}
=\sqrt{\tfrac{2}{n}}\;\cos\!\big(\Delta^\star(n')\big),
\]
hence for either column in this pair we have
\(
n\,\|q\|_\infty^2\le 2\cos^2\!\big(\Delta^\star(n')\big).
\)
If $n$ is even and $j=n/2$, the one-dimensional eigenspace yields a column with constant modulus $1/\sqrt n$, so $n\,\|q\|_\infty^2=1$.

Therefore
\begin{equation}\label{eq:M-upper}
M\big(Q^{\rm (ph)}\big)\ \le\ \max\Big\{\,\mathbf{1}_{\{2\mid n\}},\ \max_{1\le j\le \lfloor (n-1)/2\rfloor}\ 2\cos^2\!\big(\Delta^\star(n/\gcd(j,n))\big)\Big\}.
\end{equation}
Since $n'\mapsto \cos(\Delta^\star(n'))$ is strictly increasing in $n'$ in each congruence class of $n'\bmod 4$, the inner maximum is achieved when $n'$ is as large as possible, namely $n'=\max\{\,n/\gcd(j,n)\,\}=n$ (take any $j$ coprime to $n$). Using the explicit values for $\Delta^\star(n)$ from Lemma~\ref{lem:twophase} gives
\[
\max_{j}2\cos^2\!\big(\Delta^\star(n/\gcd(j,n))\big)
=
\begin{cases}
2\cos^2(\pi/n), & n\equiv 0\pmod 4,\\[1mm]
2\cos^2(\pi/2n), & n\equiv 2\pmod 4,\\[1mm]
2\cos^2(\pi/4n), & n\ \text{odd}.
\end{cases}
\]
For $n\equiv 0\pmod 4$ we have $2\cos^2(\pi/n)\ge 1$ with equality only at $n=4$; for $n\equiv 2\pmod 4$ clearly $2\cos^2(\pi/2n)>1$; for odd $n$ the $j=n/2$ column is absent. Thus the outer maximum in \eqref{eq:M-upper} equals the displayed piecewise quantity, proving
\[
M\big(Q^{\rm (ph)}\big)\ \le\
\begin{cases}
2\cos^2(\pi/n), & n\equiv 0\pmod 4,\\
2\cos^2(\pi/2n), & n\equiv 2\pmod 4,\\
2\cos^2(\pi/4n), & n\ \text{odd}.
\end{cases}
\]

In order to demonstrate optimality (the lower bound), let us consider any orthogonal matrix $Q$ obtainable from the cycle basis by independent rotations within each $E_j$ (this is what we mean by “within the cycle family”). Fix $j$ and write $n'=n/\gcd(j,n)$. By Lemma~\ref{lem:twophase}, \emph{every} orthonormal pair drawn from $E_j$ satisfies
\[
\max\{\,n\|q\|_\infty^2 : q\in E_j,\ \|q\|_2=1\,\}\ \ge\ 2\cos^2\!\big(\Delta^\star(n')\big).
\]
Indeed, fix $j$ and recall that $E_j$ is the two-dimensional eigenspace
spanned by $\{u_j,v_j\}$. By Lemma~\ref{lem:twophase}, for \emph{any}
orthonormal pair $(u,v)$ spanning $E_j$ one has
\[
\max\{\|u\|_\infty,\ \|v\|_\infty\}\ \ge\ 
\sqrt{\tfrac{2}{n}}\ \cos\!\bigl(\Delta^\star(n')\bigr),
\qquad n':=\tfrac{n}{\gcd(j,n)}.
\]
Squaring and multiplying by $n$ yields
\[
\max\{\,n\|u\|_\infty^2,\ n\|v\|_\infty^2\,\}\ \ge\
2\cos^2\!\bigl(\Delta^\star(n')\bigr).
\]
Since $u$ and $v$ are unit vectors in $E_j$, this shows the desired inequality. 
In other words, whichever orthonormal basis of $E_j$ we pick, at least
one vector in it has sup-norm squared of order $2\cos^2(\Delta^\star(n'))/n$.
Thus the displayed inequality holds for \emph{every} orthonormal pair
drawn from $E_j$. (In the special case $j=n/2$ when $n$ is even, $E_{n/2}$
is one-dimensional, and one can choose $q$ with $|q(k)|\equiv 1/\sqrt{n}$,
so $n\|q\|_\infty^2=1$; this agrees with
$2\cos^2(\Delta^\star(2))=1$.)

Hence
\[
M(Q)\ \ge\ \max\Big\{\,\mathbf{1}_{\{2\mid n\}},\ \max_{j}2\cos^2\!\big(\Delta^\star(n/\gcd(j,n))\big)\Big\}
\]
\[
=\begin{cases}
2\cos^2(\pi/n), & n\equiv 0\pmod 4,\\
2\cos^2(\pi/2n), & n\equiv 2\pmod 4,\\
2\cos^2(\pi/4n), & n\ \text{odd},
\end{cases}
\]
where the equality follows by the same monotonicity and the presence (when $2\mid n$) of the $j=n/2$ column with value $1\le$ the right-hand side. Combining the upper and lower bounds yields the asserted equality for $M(Q^{\rm (ph)})$ and proves optimality within the cycle family.

Finally, Lemma~\ref{lem:coh} with $M(Q^{\rm (ph)})$ as above gives the sufficient condition
\[
1+\sum_{j=2}^n\lambda_j>1-\tfrac{1}{M(Q^{\rm (ph)})},
\]
which is precisely the stated $\delta^{\rm (ph)}_n$.
\end{proof}

\begin{cor}[Uniform phase bound]\label{cor:uniform}
As $\delta_\ast:=\sup_{n\ge3}\delta^{\mathrm{(ph)}}_n=\tfrac12$, for every $n\ge3$ and every Sule\u{\i}manova list $(1,\lambda_2,\ldots,\lambda_n)$ with $\lambda_j\le0$,
\[
1+\sum_{j=2}^{n}\lambda_j \;>\; \delta_\ast\ (= \tfrac12)
\quad\Longrightarrow\quad
Q^{(\mathrm{ph})}\,\Lambda\,\big(Q^{(\mathrm{ph)})}\big)^{\!\top}\ \text{is symmetric, doubly stochastic, and entrywise non-negative}.
\]
\end{cor}

\section{A coherence criterion}\label{sec:coherence}

While the phase-optimised basis successfully removed the `$8 \mid n$' artefact, its coherence still approaches the maximum value of 2 as $n \to \infty$, preventing a uniform improvement on the $\delta < 1/2$ bound. To achieve a stronger result, we must look beyond the eigenbasis of the cycle graph. This section develops a general tool to guide this search. We first isolate the key property of the basis that governs the SDIEP bound, what we call its \emph{coherence}, a measure of the peak magnitude of its entries, and prove that a lower coherence directly implies a better bound $\delta$. Armed with this criterion, we then exhibit two families of `flat' bases, constructed from Hadamard matrices, which possess the ideal coherence of $M(Q)=1$. For dimensions where such matrices do exist, this allows us to prove the realisability of \emph{all} Sule\u{\i}manova spectra with a strictly positive trace sum.

\begin{lem}[Generic coherence criterion]\label{lem:coh}
Let $Q$ be any (real) orthogonal matrix with first column $q_0=\1/\sqrt n$. Define
\[
M(Q):=\sup_{1\le j\le n-1}\ \sup_{0\le k,l\le n-1}\ n\,|q_j(k)q_j(l)|
\ =\ n\max_{1\le j\le n-1}\|q_j\|_\infty^2.
\]
If $1+\sum_{j=2}^n\lambda_j>1-\tfrac1{M(Q)}$ and $\lambda_j\le0$, then $P(\Lam)=Q\Lam Q^\top$ is \emph{symmetric, doubly stochastic, and entrywise non-negative}.
\end{lem}


\begin{proof}
As in Lemma~\ref{lem:entrybound}, $P(\Lam)\1=\1$ and
\[
p_{kl}=\frac1n\Bigl(1+\sum_{j=1}^{n-1}\lambda_j\,n\,q_j(k)q_j(l)\Bigr)
\ \ge\ \frac1n\Bigl(1+M(Q)\sum_{j=1}^{n-1}\lambda_j\Bigr)
\]
since $\lambda_j\le0$ and $n\,q_j(k)q_j(l)\le M(Q)$.
\end{proof}

The cycle basis is the standard Fourier basis for the cyclic group $\mathbb{Z}_n$. To achieve the minimal coherence of $M=1$, we turn from the cycle graph to a different algebraic structure: the \emph{abelian 2-group}. An abelian 2-group is a group that is commutative and in which every element has an order that is a power of two; for a finite group, this is equivalent to its size being a power of two. The canonical example for any dimension $n=2^d$ is the elementary abelian 2-group $G = (\mathbb{Z}/2\mathbb{Z})^d$, whose elements can be viewed as binary vectors $\mathbf{x}=(x_1, \dots, x_d)$ of length $d$, and whose operation is vector addition modulo $2$ (bitwise XOR).

The natural analogue of the Fourier basis on this group is the basis of its \emph{characters}. A character is a homomorphism $\chi\colon G \to \{\pm 1\}$. For each indexing vector $\mathbf{a} \in G$, there is a corresponding character $\chi_{\mathbf{a}}: G \to \{\pm 1\}$ defined by the dot product:
\[
\chi_{\mathbf{a}}(\mathbf{x}) = (-1)^{\mathbf{a}\cdot\mathbf{x}} = (-1)^{\sum_{i=1}^d a_i x_i}.
\]
The \emph{Sylvester} (or \emph{Walsh--Hadamard}) matrix $H_n$ is constructed by indexing its rows and columns by the elements of $G$ (typically in lexicographical order) and setting the entry at row $\mathbf{a}$ and column $\mathbf{x}$ to be $\chi_{\mathbf{a}}(\mathbf{x})$.
\

\begin{rem}[Hadamard matrix orders]
It is a classical open problem, known as the \emph{Hadamard conjecture}, that Hadamard matrices exist for every order $n$ equal to $1$, $2$, or a multiple of $4$ \cite{HedayatWallis}. Therefore, a Walsh--Hadamard matrix of order $3$ does not exist.
\end{rem}

A simple recursive construction, known as the Sylvester construction, generates these matrices for all orders $n=2^d$:
\[
H_1 = \begin{pmatrix} 1 \end{pmatrix}, \quad H_2 = \begin{pmatrix} 1 & 1 \\ 1 & -1 \end{pmatrix}, \quad H_{2n} = \begin{pmatrix} H_n & H_n \\ H_n & -H_n \end{pmatrix}.
\]
For example, the matrix for $n=4=2^2$ is constructed as follows:
\[
H_4 = \begin{pmatrix} H_2 & H_2 \\ H_2 & -H_2 \end{pmatrix} = 
\left(\begin{array}{cc|cc}
1 & 1 & 1 & 1 \\
1 & -1 & 1 & -1 \\\hline
1 & 1 & -1 & -1 \\
1 & -1 & -1 & 1
\end{array}\right).
\]
By construction, the entries of $H_n$ are all $\pm 1$ and its columns are orthogonal. When normalised to $Q = H_n/\sqrt{n}$, we obtain an orthogonal matrix whose entries are all $\pm 1/\sqrt{n}$. This ``flat'' structure gives a coherence of $M(Q)=1$, providing a powerful tool for the SDIEP as formalised in the following theorem.

\begin{thm}[bases stemming from abelian $2$-groups]\label{thm:hypercube}
If $n=2^d$, let $H_n$ be the Sylvester (Walsh--Hadamard) matrix of order $n$. After sign/permutation normalisation, $Q=H_n/\sqrt n$ is orthogonal with first column $\1/\sqrt n$ and entries $\pm 1/\sqrt n$, hence $M(Q)=1$. Therefore, for \emph{every} Sule\u{\i}manova spectrum with $1+\sum_{j=2}^n\lambda_j>0$, the matrix $Q\Lam Q^\top$ is symmetric and doubly stochastic.
\end{thm}

\begin{thm}[Hadamard orders via Paley-type constructions]\label{thm:hadamard}
Whenever a (real) Hadamard matrix $H$ of order $n$ exists with first column $\1$, taking $Q=H/\sqrt n$ yields $M(Q)=1$. Hence every Sule\u{\i}manova spectrum with $1+\sum_{j=2}^n\lambda_j>0$ is realised by a symmetric doubly stochastic matrix via $Q\Lam Q^\top$.
\end{thm}

\begin{proof}[Proof of Theorems \ref{thm:hypercube} and \ref{thm:hadamard} (combined)]
Let $H$ be a real Hadamard matrix of order $n$. By multiplying columns by $\pm1$ and permuting them (operations preserving the Hadamard property), normalise so that the first column is $\1$. Set $Q:=H/\sqrt n$. Then $Q$ is real orthogonal, its first column is $q_0=\1/\sqrt n$, and all entries satisfy $q_j(k)\in\{\pm 1/\sqrt n\}$.

Hence $M(Q)=\sup_{j\ge1,k,l}n|q_j(k)q_j(l)|=1$. By Lemma~\ref{lem:coh}, whenever $1+\sum_{j=2}^n\lambda_j>1-\tfrac1{M(Q)}=0$ and $\lambda_j\le 0$, the matrix $P(\Lambda)=Q\Lambda Q^\top$ is entrywise nonnegative. Symmetry is immediate, and
\[
P(\Lambda)\,\1=Q\Lambda Q^\top\1=Q\Lambda(\sqrt n\,e_0)=\sqrt n\,Qe_0=\1
\]
so $P(\Lambda)$ is doubly stochastic. The Walsh--Hadamard case corresponds to $n=2^d$; Paley-type constructions and other Hadamard orders give further instances.
\end{proof}

\begin{rem}If $M(Q)=1$ (Hadamard orders) and $T:=1+\sum_{j=2}^n\lambda_j>\varepsilon>0$, then
$p_{k\ell}\ge T/n$ analytically. In floating point arithmetic with unit roundoff $u$,
a conservative acceptance test is $\min_{k,\ell}P_{k\ell}\ge -(cu)$ for a modest constant $c$
depending on $n$ and the BLAS implementation (we used $c=10^{-12}$ in Section~\ref{sec:numerics}).
\end{rem}


\begin{rem}
(i) The result requires the \emph{strict} condition $1+\sum\lambda_j>0$. Our lower bound on the matrix entries is $p_{kl}\ge \tfrac{1}{n}(1+\sum\lambda_j)$, which guarantees positivity under the strict inequality. The boundary case $1+\sum\lambda_j=0$ is not guaranteed to be realisable and is known to fail for certain spectra (e.g., $(1,0,\dots,0,-1)$ in odd dimensions).

(ii) Our use of Hadamard bases provides a constructive proof for \emph{all} (not necessarily sorted) Sule\u{\i}manova lists in these dimensions, provided the trace sum is positive. This complements results such as that of Johnson and Paparella \cite{JP16}, which establish existence for \emph{normalised} (sorted) Sule\u{\i}manova lists.
\end{rem}

\section{Further remarks and extensions}

\subsection*{Asymptotics (large \texorpdfstring{$n$}{n})}
For small angles, $\cos^2\theta=1-\theta^2+O(\theta^4)$, hence
\[
1-\frac{1}{2\cos^2\theta}\ =\ \frac12\ -\ \frac{\theta^2}{2}\ +\ O(\theta^4).
\]
For the canonical cycle basis, $\theta=\Delta_n=\frac{\pi}{4n}\rho(n)$ (Lemma~\ref{lem:cn}), so
\[
\delta_n\ =\
\begin{cases}
\displaystyle \frac12, & n\equiv 0\pmod{8},\\[1mm]
\displaystyle \frac12 - \frac{\pi^2}{32n^2}\rho(n)^2\ +\ O(n^{-4}), & n\not\equiv 0\pmod{8}.
\end{cases}
\]
Explicitly: $\delta_n=\frac12-\frac{\pi^2}{2n^2}+O(n^{-4})$ for $n\equiv 4\pmod{8}$; $\delta_n=\frac12-\frac{\pi^2}{8n^2}+O(n^{-4})$ for $n\equiv 2,6\pmod{8}$; and $\delta_n=\frac12-\frac{\pi^2}{32n^2}+O(n^{-4})$ for odd $n$. For the phase-optimised basis (Theorem~\ref{thm:phase-correct}),
\[
\delta_n^{\rm (ph)}=
\begin{cases}
\displaystyle \frac12-\frac{\pi^2}{2n^2}+O(n^{-4}),   & n\equiv 0\pmod{4},\\[2mm]
\displaystyle \frac12-\frac{\pi^2}{8n^2}+O(n^{-4}),  & n\equiv 2\pmod{4},\\[2mm]
\displaystyle \frac12-\frac{\pi^2}{32n^2}+O(n^{-4}), & n\ \text{odd}.
\end{cases}
\]

\subsection*{Tensor products and direct products of graphs}
If $Q^{(1)}$ and $Q^{(2)}$ are orthogonal with first columns $\1/\sqrt{n_1}$ and $\1/\sqrt{n_2}$, then $Q:=Q^{(1)}\otimes Q^{(2)}$ is orthogonal with first column $\1/\sqrt{n_1n_2}$ and
\[
M(Q)\ \le\ M\!\big(Q^{(1)}\big)\,M\!\big(Q^{(2)}\big).
\]
Thus Lemma~\ref{lem:coh} gives the sufficient sum-only condition
\[
1+\sum_{j=2}^{n_1n_2}\lambda_j\ >\ 1-\frac{1}{M(Q^{(1)})M(Q^{(2)})}
\]
for entrywise nonnegativity via $Q\Lam Q^\top$. In particular, flat ($M=1$) factors propagate.

\section{Limitations and open problems}

For the canonical cycle basis, $M(Q)=2$ exactly when $8\mid n$ (Lemma~\ref{lem:cn}), so no sum-only improvement below $\tfrac12$ is possible there. Even when $8\nmid n$, $M(Q)=2\cos^2\Delta_n=2\bigl(1-O(n^{-2})\bigr)$, so the gain is $O(n^{-2})$ (Theorem~\ref{thm:cycle}). The phase-optimised basis yields $\delta_n^{\rm (ph)}<\tfrac12$ for all $n$, but still not a single \emph{uniform} $\delta<\tfrac12$ independent of $n$.

\medskip
\noindent\emph{Open problem.} Is it possible to construct, for each $n$, a real orthogonal $Q$ with $q_0=\1/\sqrt n$ and $M(Q)\le 2-\varepsilon$ for some $\varepsilon>0$ independent of $n$. By Lemma~\ref{lem:coh}, this would give a universal $\delta<\tfrac12$ for all (not necessarily normalised) Sule\u{\i}manova lists.

\section{Numerical analysis}\label{sec:numerics}

We complement the theoretical analysis with numerical experiments. All computations were performed in Python~3.12 with NumPy~1.26.4~\cite{numpy}. Complete code, data files, and exact reproduction instructions are available in the companion software repository \texttt{sdiep-schur} at \url{https://github.com/tomaszkania/sdiep-schur-schur}. We test entrywise non-negativity with the tolerance $\min_{k,\ell} P(\Lambda)_{k\ell}\ge -10^{-12}$, which is conservative relative to round-off.

\subsection{Methodology}
For each dimension $n$ we construct orthogonal bases $Q$ with Perron column $q_0=\mathbf{1}/\sqrt n$ of four types: (i) canonical cycle (Section~\ref{sec:cycle}); (ii) phase-optimised cycle (Section~3); (iii) Walsh--Hadamard (powers of two); and (iv) Paley type~I Hadamard (orders $q+1$ with $q\equiv3\pmod4$ prime). 

Random Sule\u{\i}manova lists with prescribed trace sum $T=1+\sum_{j=2}^n\lambda_j$ are sampled as
\[
(\lambda_2,\dots,\lambda_n)=-(1-T)(X_2,\dots,X_n),\qquad (X_2,\dots,X_n)\sim\mathrm{Dirichlet}(\alpha\mathbf 1),
\]
with $\alpha=1$ (uniform on the simplex). For each draw we compute $P(\Lambda)=Q\Lambda Q^\top$ and record success if $P(\Lambda)\ge0$ entrywise. Monte-Carlo proportions are shown with Wilson~\cite{wilson1927} 95\% confidence intervals; for $N=1000$ trials the error bars are smaller than plotting symbols in our figures and are omitted for clarity.


\subsection{Coherence verification}
Figure~\ref{fig:coherence} compares the empirical coherence $M(Q)=n\max_{j\ge1}\|q_j\|_\infty^2$ to the theoretical values. The canonical cycle basis shows the mod-$8$ pattern with $M(Q)=2$ exactly at $8\mid n$; the phase-optimised basis removes this artefact and follows the piecewise form of Theorem~\ref{thm:phase-correct}. Walsh--Hadamard and Paley bases attain $M(Q)=1$ as predicted by Lemma~\ref{lem:coh}.

\begin{figure}[t]
  \centering
  \includegraphics[width=0.82\textwidth]{fig_exp1_coherence.pdf}
  \caption{Empirical coherence $M(Q)=n\max_{j\ge1}\|q_j\|_\infty^2$ versus dimension $n$ for the four basis families implemented in the numerical code. The plot is a direct numerical verification of the closed-form coherence formulae proved in Sections~\ref{sec:cycle}--\ref{sec:coherence}: the canonical cycle basis attains the extremal value $M(Q)=2$ exactly when $8\mid n$; the phase-optimised cycle basis removes this artefact and yields $M(Q)<2$ for every $n\ge3$; and Walsh--Hadamard/Paley bases attain the optimal value $M(Q)=1$. The numerical and theoretical values agree to floating-point round-off.}
  \label{fig:coherence}
\end{figure}

\begin{table}[h]
\centering
\caption{Coherence verification: relative error statistics over the tested dimensions. For the cycle and phase-optimised families this covers $n=3,\ldots,100$; for Walsh--Hadamard and Paley it covers the admissible orders in that range.}
\label{tab:coherence_summary}
\begin{tabular}{lcccr}
\toprule
Basis & Count & Mean & Std & Max \\
\midrule
Canonical cycle & 98 & $1.46\times 10^{-15}$ & $1.69\times 10^{-15}$ & $8.22\times 10^{-15}$ \\
Phase-optimised & 98 & $1.52\times 10^{-14}$ & $1.16\times 10^{-14}$ & $4.76\times 10^{-14}$ \\
Walsh--Hadamard & 5 & $8.88\times 10^{-17}$ & $1.22\times 10^{-16}$ & $2.22\times 10^{-16}$ \\
Paley type~I & 13 & $1.37\times 10^{-16}$ & $1.03\times 10^{-16}$ & $2.22\times 10^{-16}$ \\
\bottomrule
\end{tabular}
\end{table}

All discrepancies are at the level of floating-point round-off (below $5\times 10^{-14}$ in our experiments), reinforcing Lemmata~\ref{lem:cn} and~\ref{lem:twophase}.


\subsection{Success probability versus trace sum}
Figure~\ref{fig:success} plots the empirical success rate $\mathbb P[P(\Lambda)\ge0]$ as a function of the trace sum $T=1+\sum_{j\ge2}\lambda_j$ for representative sizes/bases. The panels are nearly uniformly near~$1$, i.e.\ random Dirichlet-uniform Sule\u{\i}manova lists are realised with overwhelming probability even \emph{below} the sufficient thresholds $\delta_n$ (canonical) and $\delta_n^{(\mathrm{ph})}$ (phase-optimised). This is expected: the bounds are \emph{sufficient} (worst-case) and are attained by highly structured `spiky' spectra (Proposition~\ref{prop:sharp}), whereas typical draws are far from these extremals. For Hadamard orders, success is~$1$ for all $T>0$, in agreement with Theorems~\ref{thm:hypercube}--\ref{thm:hadamard}.

\begin{figure}[p]
\centering
\includegraphics[width=0.95\textwidth]{fig_exp2_heatmap_n8.pdf}\par\medskip
\includegraphics[width=0.95\textwidth]{fig_exp2_heatmap_n12.pdf}\par\medskip
\includegraphics[width=0.95\textwidth]{fig_exp2_heatmap_n16.pdf}\par\medskip
\includegraphics[width=0.95\textwidth]{fig_exp2_heatmap_n24.pdf}
\caption{Empirical realisation probability $\mathbb{P}[P(\Lambda)\ge 0]$ as a function of the trace sum $T=1+\sum_{j\ge 2}\lambda_j$ for selected dimensions. Each block corresponds to one dimension ($n=8,12,16,24$), and within each block we compare the basis families implemented in the repository \url{https://github.com/tomaszkania/sdiep-schur-schur}. The main numerical message is a typical-case/worst-case distinction: the analytic quantities $\delta_n$ and $\delta_n^{(\mathrm{ph})}$ are worst-case sufficient thresholds, whereas for Dirichlet-sampled Sule\u{\i}manova lists the Schur template succeeds with high probability already below those values. In Hadamard orders (visible here at $n=16$), the success rate is numerically $1$ for all tested $T>0$, in agreement with Lemma~\ref{lem:coh}.}
\label{fig:success}
\end{figure}


\subsection{Sharpness at the canonical cycle threshold}
To exhibit the worst-case behaviour predicted by Proposition~\ref{prop:sharp}, we consider one-spike spectra in $n=12$: set $\lambda_{j^\ast}=t<0$ with $j^\ast=1$ (coprime to $n$) and $\lambda_j=0$ for $j\notin\{0,j^\ast\}$. By Lemma~\ref{lem:cn}, since $12\equiv4\pmod8$ we have $\rho(12)=4$ and thus $\Delta_{12}=\pi\rho(12)/(4\cdot 12)=\pi/12$. Proposition~\ref{prop:sharp} then predicts that $p_{k^\ast k^\ast}$ crosses zero at
\[
t^\ast=-\frac{1}{2\cos^2(\Delta_{12})}=-\frac{1}{2\cos^2(\pi/12)}\approx -0.535898\ldots
\]
implying $1+\sum_{j=2}^n\lambda_j=1+t^\ast\approx 0.464102$, in agreement with $\delta_{12}$ from Theorem~\ref{thm:cycle}.

Figure~\ref{fig:sharpness} shows that $\min_{k,\ell}P(\Lambda)_{k\ell}$ varies linearly in $t$ and crosses $0$ at $t\approx-0.536$, confirming the sharpness of Theorem~\ref{thm:cycle} within the cycle template.

\begin{figure}[t]
  \centering
  \includegraphics[width=0.88\textwidth]{fig_exp3_sharpness.pdf}
  \caption{Sharpness test for the canonical cycle basis at $n=12$ using one-spike spectra: $\lambda_{j^\ast}=t<0$ with $j^\ast=1$, and $\lambda_j=0$ for $j\notin\{0,j^\ast\}$. For each $t$ we evaluate the minimum entry $\min_{k,\ell}P(\Lambda)_{k\ell}$. Proposition~\ref{prop:sharp} predicts an affine dependence on $t$ and a zero crossing at
  $t^\ast=-1/(2\cos^2(\Delta_{12}))=-1/(2\cos^2(\pi/12))\approx -0.5359$.
  The numerical curve crosses $0$ at the same value (up to grid resolution and floating-point tolerance), confirming that the trace-sum bound for the canonical cycle basis is genuinely sharp within this template.}
  \label{fig:sharpness}
\end{figure}

\subsection{Computational cost}
Constructing the canonical or phase-optimised cycle basis explicitly as a dense matrix costs $O(n^2)$ floating-point operations, dominated by evaluating trigonometric functions on an $n\times n$ grid. In contrast, Walsh--Hadamard bases can be generated recursively with $O(n^2)$ memory operations and admit fast matrix--vector multiplies via the fast Walsh--Hadamard transform.

Evaluating $P(\Lambda)=Q\Lambda Q^\top$ as a dense matrix uses Level-3 BLAS operations: forming $Q\Lambda$ costs $O(n^2)$ flops (diagonal scaling), and the product $(Q\Lambda)Q^\top$ costs $O(n^3)$ flops using standard dense matrix multiplication. Representative timings are shown in Figure~\ref{fig:timing}. For moderate sizes (e.g.\ $n\le 128$) the method is inexpensive on a laptop CPU, and for repeated evaluations the structured bases (FFT/FWHT) allow significantly faster matrix--vector products without forming $P(\Lambda)$ explicitly.

\begin{figure}[t]
  \centering
  \includegraphics[width=0.88\textwidth]{fig_exp4_timing.pdf}
  \caption{Wall-clock timings for basis construction and evaluation of $P(\Lambda)=Q\Lambda Q^\top$ on the tested dimensions. Constructing the canonical and phase-optimised cycle bases scales quadratically, while explicit evaluation of $P(\Lambda)$ as a dense matrix is dominated by the matrix--matrix product $(Q\Lambda)Q^\top$ and scales cubically. The Walsh--Hadamard basis is faster to construct because of its recursive structure.}
  \label{fig:timing}
\end{figure}

\begin{rem}[Fast application without forming $P(\Lambda)$]
For many numerical tasks one does not need the full matrix $P(\Lambda)$ explicitly, but only products $P(\Lambda)x$. These can be computed as
\[
P(\Lambda)x = Q\bigl(\Lambda(Q^\top x)\bigr).
\]
Thus the cost is dominated by applying $Q$ and $Q^\top$. For Walsh--Hadamard bases this can be done by the fast Walsh--Hadamard transform in $O(n\log n)$ operations, and for Fourier/cycle-type bases by FFT-based routines after passing to the corresponding complex representation. In this sense the proposed constructions yield not only explicit realisations but also structured linear operators suitable for large-scale iterative computations.
\end{rem}



\subsection{Discussion}
The experiments support the analytic findings of the paper. Indeed,

\begin{itemize}
\item The empirical values of $M(Q)$ match the formulae exactly:
$M(Q_{\mathrm{cycle}})$ attains $2$ precisely at $8\mid n$; phase optimisation removes this
artefact and lowers $M(Q)$ across all $n$; Hadamard families achieve the optimal $M(Q)=1$.
\item For Dirichlet-uniform random Sule\u{\i}manova lists with a
given trace sum, the Schur template succeeds with overwhelming probability even
\emph{below} the sufficient thresholds $\delta_n$ and $\delta_n^{(\mathrm{ph})}$. This
is expected: the bounds are worst–case and are saturated by highly structured `spiky'
spectra, not by typical random draws.
\item In the one–spike family, the minimum entry varies
\emph{linearly} in the spike height $t$ and crosses zero at
$t^\ast=-1/(2\cos^2\Delta_n)$, exactly as predicted. Hence the cycle–basis threshold
cannot be improved by a sum–only condition within that template.
\item When a Hadamard order is available, choosing $Q$ from a
Hadamard family (Walsh or Paley) guarantees realisation for every strict
Sule\u{\i}manova list ($M(Q)=1$, hence $\delta=0$). Otherwise, the phase-optimised cycle basis
improves uniformly over the canonical choice and removes the $8\mid n$ barrier.
\end{itemize}

All code and data to reproduce Figures~\ref{fig:coherence}, \ref{fig:success}, \ref{fig:sharpness}, and~\ref{fig:timing}
are available at\newline

\url{https://github.com/tomaszkania/sdiep-schur-schur}. Additional representative values of the phase-optimised threshold are listed in Table~\ref{tab:S1} of the Appendix.

\clearpage
\appendix
\section*{Appendix: Values and visualisation of the phase-optimised bound}
\renewcommand{\thetable}{S\arabic{table}}
\setcounter{table}{0}

Selected numerical values of the phase-optimised threshold are collected in Table~\ref{tab:S1}; the full landscape up to $n=100$ is visualised in Figure~\ref{fig:supp_threshold}. In particular, Table~\ref{tab:S1} records representative values across the three congruence classes from Theorem~\ref{thm:phase-correct}.

\begin{table}[t]
\centering
\caption{Selected values of the phase-optimised threshold $\delta^{\mathrm{(ph)}}_n$.}
\label{tab:S1}
\begin{tabular}{r r @{\qquad} r r @{\qquad} r r @{\qquad} r r}
\hline
$n$ & $\delta^{\mathrm{(ph)}}_n$ & $n$ & $\delta^{\mathrm{(ph)}}_n$ & $n$ & $\delta^{\mathrm{(ph)}}_n$ & $n$ & $\delta^{\mathrm{(ph)}}_n$ \\
\hline
3  & 0.4641016151 & 4  & 0                & 5  & 0.4874571845 & 6  & 0.4641016151 \\
7  & 0.4936524043 & 8  & 0.4142135624     & 9  & 0.4961728669 & 10 & 0.4874571845 \\
12 & 0.4641016151 & 16 & 0.4802169351     & 20 & 0.4874571845 & 24 & 0.4913338099 \\
30 & 0.4986267123 & 40 & 0.4969030207     & 60 & 0.4986267123 & 100 & 0.4995061949 \\
\hline
\end{tabular}
\end{table}

\begin{figure}[t]
\centering
\begin{tikzpicture}
\begin{axis}[
  width=0.9\textwidth,
  height=7cm,
  xlabel={$n$},
  ylabel={$\,\delta^{\mathrm{(ph)}}_n$},
  xmin=3, xmax=60,
  ymin=0, ymax=0.52,
  grid=both,
  legend pos=south east,
  ytick distance=0.1,
  xtick distance=5
]
% envelopes corresponding to the three congruence classes (degrees: 180 = pi)
\addplot+[no marks, thick, domain=4:60, samples=300] {1 - 1/(2*(cos(180/x))^2)};
\addlegendentry{$n\equiv 0\pmod{4}$: $1 - \dfrac{1}{2\cos^2(\pi/n)}$}
\addplot+[no marks, thick, domain=6:60, samples=300, dashed] {1 - 1/(2*(cos(90/x))^2)};
\addlegendentry{$n\equiv 2\pmod{4}$: $1 - \dfrac{1}{2\cos^2(\pi/(2n))}$}
\addplot+[no marks, thick, domain=3:60, samples=300, densely dotted] {1 - 1/(2*(cos(45/x))^2)};
\addlegendentry{$n$ odd: $1 - \dfrac{1}{2\cos^2(\pi/(4n))}$}
% sample integer points
\addplot+[only marks] coordinates {(3,0.4641016151) (4,0) (5,0.4874571845) (6,0.4641016151) (7,0.4936524043)
(8,0.4142135624) (9,0.4961728669) (10,0.4874571845) (12,0.4641016151) (16,0.4802169351) (20,0.4874571845) (24,0.4913338099)};
%\addlegendentry{\tiny{sample $n$}}
\end{axis}
\end{tikzpicture}
\caption{Phase-optimised threshold $\delta^{\mathrm{(ph)}}_n$ by congruence class. The three curves are the analytic envelopes from Theorem~\ref{thm:phase-correct}; plotted markers indicate representative integer values from Table~\ref{tab:S1}.}\label{fig:supp_threshold}
\end{figure}

\newpage

\begin{thebibliography}{99}

\bibitem{GK20}
M.~Gnacik and T.~Kania.
\newblock Inverse problems for symmetric doubly stochastic matrices whose {Sule{\u{\i}}manova} spectra are bounded below by 1/2.
\newblock \emph{Linear Algebra and its Applications}, 592:175--187, 2020.

\bibitem{HedayatWallis}
A.~Hedayat and W.~D.~Wallis,
\newblock Hadamard matrices and their applications,
\emph{Ann. Statist.} \textbf{6} (1978), 1184--1238.

\bibitem{Horadam}
K.~J.~Horadam,
\newblock \emph{Hadamard Matrices and Their Applications}, Princeton University Press, 2007.

\bibitem{JP16}
C.~R. Johnson and P.~Paparella.
\newblock Perron spectratopes and the real nonnegative inverse eigenvalue problem.
\newblock \emph{Linear Algebra and its Applications}, 493:281--300, 2016.

\bibitem{numpy}
C.~R. Harris \emph{et al.}, 
\newblock Array programming with NumPy, 
\emph{Nature} \textbf{585} (2020), 357--362.

\bibitem{matplotlib}
J.~D. Hunter, 
\newblock Matplotlib: A 2D graphics environment, 
\emph{Computing in Science \& Engineering} \textbf{9} (2007), no.~3, 90--95.

\bibitem{Paparella16}
P.~Paparella.
\newblock Realizing {Suleimanova}-type spectra via permutative matrices.
\newblock \emph{Electronic Journal of Linear Algebra}, 31:306--312, 2016.

\bibitem{Paley}
R.~E.~A.~C.~Paley,
\newblock On orthogonal matrices, \emph{J. Math. Phys.} \textbf{12} (1933), 311--320.

\bibitem{Perfect}
H.~Perfect,
\newblock Methods of constructing certain stochastic matrices~I,
\emph{Duke Math. J.} \textbf{20} (1953), 395--404.

\bibitem{PM65}
H.~Perfect and L.~Mirsky.
\newblock Spectral properties of doubly-stochastic matrices.
\newblock \emph{Monatshefte f{\"u}r Mathematik}, 69:35--57, 1965.

\bibitem{Sarkhoni25}
M.~R. Sarkhoni, H.~M. Kermani, and A.~Rivaz.
\newblock Symmetric doubly stochastic inverse eigenvalue problem for odd sizes.
\newblock \emph{Linear Algebra and its Applications}, 708:594--607, 2025.

\bibitem{Suleimanova65}
H.~R. Sule{\u{\i}}manova.
\newblock The question of a necessary and sufficient condition for the existence of a stochastic matrix with prescribed characteristic numbers.
\newblock \emph{Trudy Vsesojuz. Zaočn. Ènerget. Inst. Vyp.}, 28:33--49, 1965. (in Russian).

\bibitem{Walsh}
J.~L.~Walsh,
\newblock A closed set of normal orthogonal functions, \emph{Amer. J. Math.} \textbf{45} (1923), 5--24.

\bibitem{wilson1927}
E.~B. Wilson, 
\newblock Probable inference, the law of succession, and statistical inference, 
\emph{Journal of the American Statistical Association} \textbf{22} (1927), no.~158, 209--212.

\end{thebibliography}

\end{document}